\section{Lecture 5, Part 0: The Central Limit Theorem}

\begin{frame}
\vfill
\centering
{\color{RLBlue}\large Lecture 5}

\vspace{0.3em}
{\color{RLBlue}\Large\bfseries Part 0}

\vspace{0.6em}
{\Large The Central Limit Theorem}
\vfill
\end{frame}

\begin{frame}{Learning Goals}
\begin{itemize}
  \item Treat the sample mean $\bar{x}$ as a random variable, and describe its sampling distribution.
  \pause
  \item Use the standard error $\sigma/\sqrt{n}$, and explain why four times the data buys only half the error.
  \pause
  \item State the central limit theorem with its conditions, and say what goes wrong when they fail.
\end{itemize}
\end{frame}

% ─── The sampling distribution ──────────────────────────────────────────────
\begin{frame}{Why Do We Need a Sampling Distribution?}
\begin{itemize}
  \item \textbf{Why:} every average we have ever computed came from \emph{one} sample. Collect a different sample and the average changes. So $\bar{x}$ is not a fixed number -- it is a random variable with a distribution of its own.
  \pause
  \item \textbf{Numerical example:} the coffee shop sells $0$, $1$ or $2$ pastries per order with probabilities $0.7, 0.2, 0.1$ (Lecture 3), so $\mu = 0.4$ and $\sigma^2 = 0.44$.
  \pause
  \item Serve two customers. You might see $0$ and $0$ (average $0$), or $1$ and $2$ (average $1.5$). Same shop, same day, different answer.
  \pause
  \item Until we know how much $\bar{x}$ moves around, a single measured average tells us nothing about how much to trust it.
\end{itemize}
\end{frame}

\begin{frame}{The Sampling Distribution of the Mean}
\begin{rldefinition}
The \textbf{sampling distribution} of $\bar{x}$ is the distribution of the averages you would get by repeating the whole sampling exercise again and again. \emph{Everyday:} ask $30$ customers their order size, write down the average, send $30$ fresh customers, repeat -- the histogram of those averages is the sampling distribution.
\[
E[\bar{x}] = \mu, \qquad \mathrm{Var}(\bar{x}) = \frac{\sigma^2}{n} .
\]
\end{rldefinition}
\pause
\begin{itemize}
  \item The first says $\bar{x}$ is centred on the truth -- it is \textbf{unbiased}. The second says its spread shrinks as the sample grows.
  \pause
  \item Neither requires the original distribution to have any particular shape. Both are exact for every $n$.
\end{itemize}
\end{frame}

\begin{frame}{Worked Example: Two Orders, Five Possible Averages}
\begin{itemize}
  \item Serve $n = 2$ customers. Every pair of orders and its average:
\end{itemize}
\centering
\small
\begin{tabular}{l|ccccc}
Average $\bar{x}$    & $0$    & $0.5$  & $1$    & $1.5$  & $2$    \\
\hline
Probability          & $0.49$ & $0.28$ & $0.18$ & $0.04$ & $0.01$ \\
\end{tabular}

\vspace{0.4em}
\raggedright
\begin{itemize}
  \item[] \pause
  \item $E[\bar{x}] = 0(0.49) + 0.5(0.28) + 1(0.18) + 1.5(0.04) + 2(0.01) = 0.40 = \mu$. \pause
  \item $E[\bar{x}^2] = 0.25(0.28) + 1(0.18) + 2.25(0.04) + 4(0.01) = 0.38$, so $\mathrm{Var}(\bar{x}) = 0.38 - 0.4^2 = 0.22$. \pause
  \item And $\sigma^2/n = 0.44/2 = 0.22$ -- the formula checked by hand, with no appeal to large samples.
\end{itemize}
\end{frame}

% ─── Standard error ─────────────────────────────────────────────────────────
\begin{frame}{Why Do We Need the Standard Error?}
\begin{rldefinition}
The \textbf{standard error} is the standard deviation of $\bar{x}$ -- how far a computed average typically sits from the truth:
\[
\mathrm{SE} = \frac{\sigma}{\sqrt{n}} .
\]
\emph{Not the same as $\sigma$.} $\sigma$ says how much one \emph{customer} varies; SE says how much an \emph{average of $n$ customers} varies.
\end{rldefinition}
\pause
\begin{itemize}
  \item \textbf{Coffee shop:} $\sigma = \sqrt{0.44} = 0.663$. One order is wildly unpredictable, but the average of $30$ has $\mathrm{SE} = 0.663/\sqrt{30} = 0.121$ -- five times steadier.
  \pause
  \item The $\sqrt{n}$ is the whole story: it is what makes measuring more worthwhile, and what makes it expensive.
\end{itemize}
\end{frame}

\begin{frame}{How Fast Does the Average Settle?}
\begin{columns}[T]
\column{0.50\textwidth}
\small
\begin{itemize}
  \item Standard error against sample size, for the coffee shop's $\sigma = 0.663$.
  \pause
  \item $n = 10 \Rightarrow 0.210$; \; $n = 40 \Rightarrow 0.105$; \; $n = 160 \Rightarrow 0.052$.
  \pause
  \item \textbf{Read it:} each step \emph{quadruples} the data and \emph{halves} the error. There is no sample size at which the curve becomes steep again.
  \pause
  \item \textbf{Why this matters:} accuracy is bought at quadratic cost. Going from $1\%$ error to $0.1\%$ needs $100$ times the data -- the reason careful measurement is expensive everywhere, from opinion polls to model evaluation.
\end{itemize}
\column{0.48\textwidth}
\centering
\includegraphics[width=\linewidth]{figures/clt_mle/clt_mle_standard_error.pdf}
\end{columns}
\end{frame}

% ─── The CLT itself ─────────────────────────────────────────────────────────
\begin{frame}{The Central Limit Theorem}
\begin{rldefinition}
Take independent draws $X_1, \dots, X_n$ from \emph{any} distribution with mean $\mu$ and finite variance $\sigma^2$. As $n$ grows, the sampling distribution of $\bar{x}$ approaches a normal:
\[
\bar{x} \;\approx\; \mathcal{N}\!\left(\mu, \; \frac{\sigma^2}{n}\right)
\qquad\text{equivalently}\qquad
\frac{\bar{x} - \mu}{\sigma/\sqrt{n}} \;\approx\; \mathcal{N}(0, 1).
\]
\emph{The striking part:} the original shape does not matter -- coin flips, dice, waiting times and pastries per order all have bell-shaped averages.
\end{rldefinition}
\pause
\begin{itemize}
  \item Lecture 3's law of large numbers said \emph{where} $\bar{x}$ lands; the CLT says what its \emph{whole distribution} looks like on the way.
  \pause
  \item \textbf{Three conditions:} independent draws, from the same distribution, with finite variance.
\end{itemize}
\end{frame}

\begin{frame}{Real Event: Three Spikes Become a Bell}
\begin{columns}[T]
\column{0.50\textwidth}
\small
\begin{itemize}
  \item Bars are the \emph{exact} distribution of $\bar{x}$, computed by combining every possible set of orders. The curve is $\mathcal{N}(0.4, 0.44/n)$.
  \pause
  \item $n = 1$ is three bars and nothing like a bell. $n = 2$ is lopsided. By $n = 5$ the shape is recognisable, and $n = 30$ tracks the curve closely.
  \pause
  \item Check it numerically: $P(\bar{x} \text{ within } 1\,\mathrm{SE})$ runs $0.90 \to 0.77 \to 0.67 \to 0.66$, closing on the $0.68$ the Gaussian predicts (Lecture 4).
  \pause
  \item \textbf{Nothing was simulated.} The bell is arithmetic, not a curve drawn over a claim.
\end{itemize}
\column{0.48\textwidth}
\centering
\includegraphics[width=\linewidth]{figures/clt_mle/clt_mle_sampling_distribution.pdf}
\end{columns}
\end{frame}

\begin{frame}{Real Event: When the CLT Is Slow}
\begin{columns}[T]
\column{0.50\textwidth}
\small
\begin{itemize}
  \item Same $n = 30$, two populations. Below, one order in twenty is an office run of $40$ coffees.
  \pause
  \item The office-run averages come in \emph{clumps} -- one per number of big orders that happened to land in the sample. The normal curve is hopeless.
  \pause
  \item Skewness measures lopsidedness -- $0$ for a symmetric shape. The skewness of $\bar{x}$ is $\mathrm{skew}(X)/\sqrt{n}$: $0.26$ above, $0.75$ below.
  \pause
  \item \textbf{So ``$n \ge 30$'' is a rule of thumb, not a theorem.} A rare huge value is exactly what breaks it.
\end{itemize}
\column{0.48\textwidth}
\centering
\includegraphics[width=\linewidth]{figures/clt_mle/clt_mle_heavy_tail.pdf}
\end{columns}
\end{frame}

\begin{frame}{Why This Matters: Mini-Batch Gradients}
\begin{itemize}
  \item \textbf{The promise, cashed:} training a network computes a gradient. The honest one averages over all $N$ training examples; that is far too slow, so a \textbf{mini-batch} of $B$ examples is averaged instead.
  \pause
  \item A mini-batch gradient is a sample mean. So it is \emph{unbiased} -- centred on the full-data gradient -- and its noise has standard error $\sigma_g/\sqrt{B}$. Training works at all because of this slide.
  \pause
  \item \textbf{Repeat this one:} \emph{quadrupling the batch size halves the gradient noise.} That is why batch $2048$ is not $32$ times better than batch $64$ -- it is only $\sqrt{32} \approx 5.7$ times less noisy.
  \pause
  \item \textbf{Batch normalization} rests on it too: it normalises by a mean and variance estimated from $B$ examples. At $B = 2$ or $4$ those estimates are themselves noisy -- one reason layer and group normalization exist (Module 2).
\end{itemize}
\end{frame}
